% ------------------------------------------------------------------------------
% Abstrak (Bahasa Indonesia)
% ------------------------------------------------------------------------------
\begin{center}
\textbf{\large ABSTRAK}\\
\end{center}
\vspace{1cm}

\begingroup
\fontsize{11pt}{13pt}\selectfont % Mengatur ukuran font jadi 11pt, spasi 1 (13pt leading)

Abstrak ditulis dengan format font Time New Roman size 11 spasi 1 dalam bahasa Indonesia. Isinya merupakan uraian singkat tetapi lengkap mengenai permasalahan, tujuan, metode penelitian, hasil, dan kesimpulan. Tujuan penelitian disarikan dari tujuan penelitian pada pendahuluan, metode penelitian disarikan dari proses penelitian, dan hasil penelitian ditarik dari kesimpulan dengan nilai hasil secara kualitatif dan kuantitatif. Abstrak atau abstract terdiri atas satu alinea dan satu spasi maksimal 200 kata. Tidak boleh ada sitasi pada abstrak. Pada abstrak tidak menggunakan penamaan, simbol atau istilah yang teknis, misalnya minsup untuk menyatakan nilai support minimal.

\vspace{0.5cm}
\noindent \textbf{Kata Kunci:} merupakan kata-kata kunci yang menjelaskan isi tulisan, biasanya dapat diambil dari judul dan abstrak. Maksimal enam buah dan ditulis dengan huruf kecil. Contoh: abstrak, bahasa Indonesia, kualitatif, kuantitatif
\endgroup